% generated by nfv/paper/make_compare_tables.py -- do not edit by hand
\begin{table}[t]\centering\footnotesize
\caption{Every entry on which a verdict and a ground-truth label disagree because the specification omits a precondition. \emph{exposed by} is the task whose claim revealed it. \emph{omitted} is the fact that, had it been stated, would have made the verdict and the label agree; all \anprecdisptotal{} rows name one, and only \texttt{116} has it written in the docstring already. The reference of a pair keeps its label by policy, so the omission changes a label only for \texttt{116}, which is a mutant. Every row is executed rather than asserted.}
\label{tab:flipped}
\begin{tabular}{lllp{0.13\textwidth}p{0.11\textwidth}p{0.28\textwidth}}
\toprule
\# & function & exposed by & omitted & effect & why \\
\midrule
\texttt{11} & \texttt{string\_xor} & refutation & $|a| = |b|$ & label kept & \footnotesize postcondition asserts $|rv|=|a|=|b|$ but nothing requires equal-length operands, so no implementation satisfies it when they differ \\
\texttt{20} & \texttt{find\_closest\_elements} & refutation & no \texttt{NaN} in \texttt{numbers} & label kept & \footnotesize witness is a pair of equal elements with $\neg(x \le x)$, admitted because \texttt{PyFloat} is abstract and realisable in Python as \texttt{NaN}, so no pair satisfies $x \le y$ \\
\texttt{49} & \texttt{modp} & refutation & $p > 0$ & label kept & \footnotesize witness is a \emph{negative} modulus (\texttt{n == 0}, \texttt{p == -1}); the docstring (``return $2^n$ modulo $p$'') fixes no sign, and under $p > 0$ the reference holds on every sampled input \\
\texttt{64} & \texttt{vowels\_count} & refutation & \texttt{s} is ASCII & label kept & \footnotesize specification counts vowels in \texttt{s.lower()} while the program tests the un-lowercased character, and \texttt{\textbackslash u0130} lowercases to two code points \\
\texttt{79} & \texttt{decimal\_to\_binary} & refutation & $decimal \ge 0$ & label kept & \footnotesize witness is a negative input: \texttt{bin(-3)[2:]} is \texttt{'b11'}, so the result carries a non-binary character and fails the \texttt{\textasciicircum db[01]+db\$} match \\
\texttt{111} & \texttt{histogram} & refutation & tokens are single letters & label kept & \footnotesize postcondition scores a key with \texttt{test.count(key)}, which counts \emph{substrings} while the program counts whitespace-separated words; the two agree only when every token is one character \\
\texttt{116} & \texttt{sort\_array} & verification & \texttt{arr} is non-negative \emph{(stated)} & label $\to$ \emph{correct} & \footnotesize the two implementations differ in the sort key only for negative values, which the docstring's ``non-negative integers'' excludes \\
\bottomrule
\end{tabular}
\end{table}
